\documentclass[11pt,a4paper]{article}
\usepackage[margin=1in]{geometry}
\usepackage{times}
\usepackage{latexsym}
\usepackage{amsmath}
\usepackage{amssymb}
\usepackage{booktabs}
\usepackage{multirow}
\usepackage{graphicx}
\usepackage{xcolor}
\usepackage{subcaption}
\usepackage[hidelinks]{hyperref}
\usepackage{natbib}
\renewcommand{\citep}{\cite}
\renewcommand{\citet}{\cite}

\title{Evidence for Architecture-Modulated Cross-Axis Signals\\in Reasoning Inference}

\author{Anonymous}

\begin{document}
\maketitle

% ===========================================================================
\begin{abstract}
Reasoning-capable language models produce long thinking traces with substantial token-level redundancy, opening two complementary efficiency axes: KV cache eviction and MLP sparsification.
We investigate whether a single, lightweight signal---EMA-based KV similarity---can coordinate both axes simultaneously (\emph{cross-axis redundancy hypothesis}).
Experiments across four model--training configurations arranged in a $2\times 2$ factorial design reveal that the cross-axis signal is \emph{architecture-modulated}: Llama-3.1-8B (GQA, 4:1 group ratio) exhibits strong correlation ($r{=}0.47$) and positive perturbation gain ($\Delta_\text{pert}{=}{+}6.6\%$), while Qwen2-based models show markedly weaker signals ($r{\le}0.25$) regardless of training paradigm.
Architecture is the dominant factor ($\Delta r{\approx}0.18$), with training paradigm contributing a secondary additive effect ($\Delta r{\approx}0.10$).
A third GQA model (Mistral-7B) confirms the positive direction ($r{=}0.31$, $\Delta_\text{pert}{=}{+}2.8\%$) but with substantially weaker magnitude, revealing within-family variation.
As a standalone contribution, our O(1) KV eviction method (CARC-KV) achieves near-parity perplexity with O($n$) baselines (H2O, SnapKV) at 70\% KV budget while providing up to 26$\times$ eviction speedup at 32K sequence length.
Per-layer decomposition of the cross-axis signal shows that mid-layer transformer blocks (layers 8--23) carry the strongest signal, with a peak per-layer $r{=}0.61$ at layer 18 in Llama.
\end{abstract}

% ===========================================================================
\section{Introduction}
\label{sec:intro}

Large reasoning models such as DeepSeek-R1 \cite{deepseekai2025deepseekr1} and QwQ \cite{qwq32b} generate extended thinking traces---often exceeding 10K tokens---to solve mathematical and scientific problems.
These traces exhibit substantial token-level redundancy: exploratory steps, self-corrections, and restated intermediate results collectively inflate the KV cache and activate MLP neurons that may contribute little to the final answer.

Two established efficiency axes target this redundancy independently.
\emph{KV cache eviction} methods (H2O \cite{zhang2024h2o}, SnapKV \cite{li2024snapkv}, StreamingLLM \cite{xiao2024streamingllm}) reduce memory by pruning less-important KV entries.
\emph{MLP sparsification} methods (Deja Vu \cite{liu2023dejavu}, CATS \cite{lee2024cats}) skip inactive neurons to reduce compute.
These two axes share a common root cause---token-level redundancy---yet are treated as orthogonal in the existing literature.

We ask: \emph{can a single redundancy signal drive both KV eviction and MLP sparsity allocation simultaneously?}
Concretely, we propose the \textbf{cross-axis redundancy hypothesis}: tokens whose KV representations are highly similar to a running EMA anchor are ``redundant'' in the attention sense, carry low surprisal, and therefore tolerate aggressive MLP sparsification.
If true, a single O(1) similarity computation could coordinate both efficiency axes without requiring separate attention-score or activation-magnitude passes.

Our experiments reveal that this hypothesis is \emph{not universally true or false} but rather \textbf{architecture-modulated}.
On Llama-3.1-8B with standard GQA attention (4:1 query-to-KV group ratio), the cross-axis signal is robust: the Pearson correlation between token-level KV similarity and MLP sparsity tolerance is $r{=}0.47$ ($p{<}0.001$), and KV-similarity-guided sparsity outperforms uniform sparsity by $\Delta_\text{pert}{=}{+}6.6\%$ at 55\% sparsity.
On Qwen2-based architectures with a 7:1 group ratio, the signal weakens dramatically ($r{\le}0.25$), and on the reasoning-distilled variant (DeepSeek-R1-Distill-Qwen-7B), the perturbation effect inverts ($\Delta_\text{pert}{=}{-}2.2\%$).

A factorial analysis ($2\times 2$: architecture $\times$ training paradigm) quantifies the contributions: architecture accounts for $\Delta r{\approx}0.18$ (main effect), training paradigm for $\Delta r{\approx}0.10$ (secondary), and the two factors are approximately additive.
A third GQA model, Mistral-7B-Instruct-v0.3, confirms the positive direction of the signal ($r{=}0.31$) but at reduced magnitude, demonstrating that even within the GQA family, model-specific factors modulate signal strength.

Beyond the cross-axis investigation, we contribute:
\begin{itemize}
    \item \textbf{CARC-KV}, an O(1) KV eviction method based on EMA anchor similarity that achieves perplexity parity with O($n$) methods at 70\% KV budget and provides 3--26$\times$ eviction speedup across sequence lengths 4K--32K.
    \item \textbf{Per-layer signal decomposition} showing that the cross-axis correlation concentrates in mid-layer transformer blocks (layers 8--23) with a peak $r{=}0.61$, providing mechanistic insight into where architecture differences manifest.
    \item \textbf{Per-head analysis} ruling out KV head count as a confound: the per-head signal ratio (3.6$\times$) exceeds the aggregate ratio (3.1$\times$).
\end{itemize}

% ===========================================================================
\section{Related Work}
\label{sec:related}

\paragraph{KV Cache Compression.}
H2O \cite{zhang2024h2o} maintains a ``heavy hitter'' set based on cumulative attention scores, requiring O($n$) updates per decoding step.
SnapKV \cite{li2024snapkv} uses an observation window to identify important KV entries, also O($n$).
StreamingLLM \cite{xiao2024streamingllm} retains only attention sinks and recent tokens, achieving O(1) cost but discarding semantically important middle tokens.
Token Filtering \cite{cao2024tokenfiltering} proposed redundancy-based KV selection but did not connect to MLP sparsity.
Our CARC-KV method uses EMA anchor similarity for O(1) eviction decisions while maintaining quality competitive with O($n$) methods.

\paragraph{MLP Sparsification.}
Deja Vu \cite{liu2023dejavu} predicts MLP sparsity patterns from input activations.
CATS \cite{lee2024cats} and PowerInfer \cite{song2024powerinfer} exploit activation sparsity for inference acceleration.
CLADA \cite{song2024clada} connects attention patterns to activation sparsity, providing indirect evidence for cross-axis coupling.
These methods treat MLP sparsification as an independent axis; we investigate whether KV-derived signals can inform sparsity allocation.

\paragraph{Efficient Reasoning.}
Recent work on reasoning-model efficiency includes thinking-token compression \cite{chen2024thinkingcompression}, early exit strategies \cite{wang2024earlyexit}, and budget allocation \cite{muennighoff2025s1}.
Our work is complementary: we focus on within-generation efficiency (per-token KV and MLP savings) rather than sequence-level truncation.

\paragraph{Attention Architecture Variants.}
Multi-Query Attention (MQA) \cite{shazeer2019mqa}, Grouped-Query Attention (GQA) \cite{ainslie2023gqa}, and Multi-head Latent Attention (MLA) \cite{liu2024deepseekv2} represent a spectrum of KV compression strategies.
We show that these architectural choices have downstream effects on the viability of semantic similarity-based efficiency signals.

% ===========================================================================
\section{Method: Cross-Axis Redundancy Coordination (CARC)}
\label{sec:method}

\subsection{EMA Anchor Similarity}

For each KV head $h$ and layer $l$, we maintain an exponential moving average (EMA) anchor $\mathbf{a}^{(l,h)}_t$ over the key vectors:
\begin{equation}
    \mathbf{a}^{(l,h)}_t = \gamma \cdot \mathbf{a}^{(l,h)}_{t-1} + (1 - \gamma) \cdot \mathbf{k}^{(l,h)}_t
\end{equation}
where $\gamma{=}0.99$ is the decay factor and $\mathbf{k}^{(l,h)}_t$ is the key vector at step $t$.
The \emph{redundancy score} for token $t$ is the cosine similarity between $\mathbf{k}^{(l,h)}_t$ and the anchor:
\begin{equation}
    s_t = \frac{1}{|\mathcal{H}|} \sum_{h \in \mathcal{H}} \cos(\mathbf{k}^{(l,h)}_t, \mathbf{a}^{(l,h)}_t)
\end{equation}
averaged across heads $\mathcal{H}$ and aggregated across layers by averaging.
Each update and similarity computation is O(1) per token per head---independent of sequence length.

\subsection{CARC-KV: O(1) KV Eviction}

When the KV cache exceeds a budget of $B$ entries, the token with the highest redundancy score (most similar to its EMA anchor) is evicted.
Unlike H2O, which requires maintaining and sorting cumulative attention scores over all $n$ cached tokens (O($n$) per step), CARC-KV eviction is O(1) amortized: we maintain a min-heap over redundancy scores and update only the current token's entry.

\subsection{Cross-Axis MLP Sparsity Allocation}

The cross-axis hypothesis posits that tokens with high KV redundancy scores tolerate more aggressive MLP sparsification.
We implement a \emph{two-tier} sparsity scheme:
\begin{equation}
    \text{sparsity}(t) = \begin{cases}
        \alpha_{\text{high}} & \text{if } s_t > \tau \\
        \alpha_{\text{low}}  & \text{otherwise}
    \end{cases}
\end{equation}
where $\alpha_{\text{high}}$ and $\alpha_{\text{low}}$ are the sparsity levels for redundant and novel tokens respectively, and $\tau$ is determined by the EMA anchor similarity distribution.

\subsection{Evaluation Metrics}

We define two primary metrics for evaluating the cross-axis signal:

\paragraph{Pearson correlation ($r$).}
Token-level Pearson correlation between KV redundancy score $s_t$ and MLP sparsity tolerance, measured as the maximum sparsity at which token $t$'s output remains within $\epsilon$ of the dense output.

\paragraph{Perturbation gain ($\Delta_\text{pert}$).}
The accuracy difference between KV-similarity-guided and uniform MLP sparsification at matched average sparsity level:
\begin{equation}
    \Delta_\text{pert} = \text{Acc}_{\text{guided}} - \text{Acc}_{\text{uniform}}
\end{equation}

% ===========================================================================
\section{Experimental Setup}
\label{sec:setup}

\subsection{Models}

We evaluate five models spanning two architecture families and two training paradigms:

\begin{itemize}
    \item \textbf{Llama-3.1-8B-Instruct} (GQA, 8 KV heads, 4:1 group ratio, instruction-tuned)
    \item \textbf{DeepSeek-R1-Distill-Llama-8B} (GQA, same Llama architecture, reasoning-distilled)
    \item \textbf{Qwen2-7B-Instruct} (GQA, 4 KV heads, 7:1 group ratio, instruction-tuned)
    \item \textbf{DeepSeek-R1-Distill-Qwen-7B} (GQA, same Qwen2 architecture, reasoning-distilled)
    \item \textbf{Mistral-7B-Instruct-v0.3} (GQA, 8 KV heads, sliding window attention)
\end{itemize}

The first four models form a $2 \times 2$ factorial design (architecture $\times$ training paradigm), enabling isolation of each factor's contribution.
Mistral-7B serves as an independent GQA validation point.

\subsection{Evaluation Protocol}

For cross-axis experiments, we use 20 AIME competition problems as reasoning benchmarks, measuring per-problem and aggregate correlation ($r$) and perturbation gain ($\Delta_\text{pert}$) at MLP sparsity levels from 40\% to 70\%.
EMA decay is fixed at $\gamma{=}0.99$ throughout.
For KV eviction benchmarks, we report perplexity on reasoning traces and AIME accuracy at KV budgets of 30\%, 50\%, and 70\%.

% ===========================================================================
\section{Results}
\label{sec:results}

\subsection{Cross-Axis Signal is Architecture-Modulated}

Table~\ref{tab:confound} presents the $2\times 2$ factorial results.
Architecture is the dominant factor: switching from Llama to Qwen2 architecture reduces $r$ by approximately 0.18 on average across training paradigms ($\bar{r}_{\text{Llama}}{=}0.425$ vs.\ $\bar{r}_{\text{Qwen2}}{=}0.20$).
Training paradigm contributes a secondary effect of $\Delta r{\approx}0.10$ (instruction $>$ reasoning-distill), and the two factors are approximately additive: the predicted value for Qwen2+Reasoning ($0.47 - 0.18 - 0.10 = 0.19$) is close to the observed $r{=}0.15$.

\begin{table}[t]
\centering
\small
\begin{tabular}{llcc}
\toprule
\textbf{Architecture} & \textbf{Training} & \textbf{$r$} & \textbf{$\Delta_\text{pert}$} \\
\midrule
Llama-GQA (4:1) & Instruction      & 0.47 & +6.6\% \\
Llama-GQA (4:1) & Reasoning-distill & 0.38 & +4.0\% \\
\midrule
Qwen2-GQA (7:1) & Instruction      & 0.25 & +1.5\% \\
Qwen2-GQA (7:1) & Reasoning-distill & 0.15 & $-$2.2\% \\
\bottomrule
\end{tabular}
\caption{$2\times 2$ factorial design isolating architecture and training effects on the cross-axis signal at 55\% MLP sparsity. Architecture is the dominant factor ($\Delta r{\approx}0.18$), training paradigm the secondary factor ($\Delta r{\approx}0.10$), with near-additive interaction.}
\label{tab:confound}
\end{table}

The perturbation gain $\Delta_\text{pert}$ follows a consistent pattern: positive and practically meaningful on Llama ($+6.6\%$ and $+4.0\%$), marginal on Qwen2-Instruct ($+1.5\%$), and inverted on Qwen2+Reasoning ($-2.2\%$).
This demonstrates that the cross-axis signal not only varies in strength but can become counterproductive when the underlying KV representations lack semantic fidelity in the similarity space.

\subsection{Within-GQA Variation: Mistral-7B}

Mistral-7B-Instruct-v0.3 provides a second GQA data point with the same KV head count (8) as Llama.
Table~\ref{tab:gqa_comparison} shows that the signal direction is consistent ($r{=}0.31$, $\Delta_\text{pert}{=}{+}2.8\%$) but magnitude is substantially reduced.
Per-problem consistency also decreases: 4 of 20 problems show $r{<}0.2$ on Mistral, compared to 0 on Llama.

\begin{table}[t]
\centering
\small
\begin{tabular}{lcccc}
\toprule
\textbf{Model} & \textbf{KV Heads} & \textbf{$r$} & \textbf{$\Delta_\text{pert}$} & \textbf{$r{<}0.2$ Problems} \\
\midrule
Llama-3.1-8B   & 8 & 0.47 & +6.6\% & 0/20 \\
Mistral-7B     & 8 & 0.31 & +2.8\% & 4/20 \\
DS-R1-Qwen-7B  & 4 & 0.15 & $-$2.2\% & 12/20 \\
\bottomrule
\end{tabular}
\caption{Cross-axis signal across GQA models. Mistral confirms the positive direction but at reduced magnitude, revealing within-family variation. DS-R1-Qwen-7B uses the Qwen2 architecture with a 7:1 group ratio.}
\label{tab:gqa_comparison}
\end{table}

The Mistral result prevents a simple ``GQA works'' narrative.
Possible explanations for the attenuation include Mistral's sliding window attention (which limits long-range KV context in some layers) and differences in training data and objectives.
The practical significance of $\Delta_\text{pert}{=}{+}2.8\%$ at 55\% sparsity warrants honest assessment: this translates to modest downstream gains that may not justify the additional complexity of cross-axis coordination for all deployment scenarios (see \S\ref{sec:downstream}).

\subsection{Per-Layer Signal Decomposition}

Table~\ref{tab:perlayer} presents a per-layer decomposition of the cross-axis correlation, binned into 8-layer groups.
On Llama, the signal peaks in mid-layer blocks (layers 8--23, $r{=}0.52$--$0.55$) with a single-layer peak of $r{=}0.61$ at layer 18.
The Qwen2-architecture model (DeepSeek-R1-Distill-Qwen-7B) shows uniformly weak signal across all layer groups ($r{\le}0.18$), with the largest gap appearing in mid-layers ($\Delta r{=}0.34$--$0.39$).

\begin{table}[t]
\centering
\small
\begin{tabular}{lccc}
\toprule
\textbf{Layer Range} & \textbf{Llama $r$} & \textbf{Qwen2 $r$} & \textbf{Gap} \\
\midrule
0--7   & 0.28 & 0.15 & 0.13 \\
8--15  & 0.52 & 0.18 & 0.34 \\
16--23 & 0.55 & 0.16 & 0.39 \\
24--31 & 0.44 & 0.12 & 0.32 \\
\bottomrule
\end{tabular}
\caption{Per-layer cross-axis correlation comparison (8-layer bins). The architecture gap is largest in mid-layers where the signal is strongest on Llama.}
\label{tab:perlayer}
\end{table}

This layer-wise pattern aligns with the functional specialization of transformer blocks: early layers build low-level representations (weak semantic content in KV similarity for both architectures), middle layers perform core reasoning transformations (strong coupling between KV redundancy and MLP activation patterns on Llama), and final layers transition toward output distributions (partial decoupling).
Twenty of Llama's 32 layers exhibit $r{>}0.4$, suggesting the signal is broadly distributed rather than concentrated in a few layers.

\subsection{Per-Head Analysis Rules Out Head-Count Confound}

A natural concern is whether the aggregate signal difference simply reflects KV head count (Llama: 8, Qwen2: 4).
Table~\ref{tab:perhead} shows that even on a per-head basis, Llama's signal is $3.6\times$ stronger than Qwen2's (mean per-head $r$: 0.40 vs.\ 0.11), exceeding the aggregate ratio of $3.1\times$.
This rules out head count as the primary driver and implicates architectural factors such as the query-to-KV group ratio (4:1 vs.\ 7:1), which determines how many query heads each KV head must serve.

\begin{table}[t]
\centering
\small
\begin{tabular}{lcccc}
\toprule
\textbf{Model} & \textbf{Heads} & \textbf{Per-Head $\bar{r}$} & \textbf{Range} & \textbf{Agg.\ $r$} \\
\midrule
Llama-3.1-8B     & 8 & 0.40 & [0.28, 0.52] & 0.47 \\
DS-R1-Qwen-7B    & 4 & 0.11 & [0.04, 0.19] & 0.15 \\
\bottomrule
\end{tabular}
\caption{Per-head cross-axis signal comparison. The per-head ratio ($3.6\times$) exceeds the aggregate ratio ($3.1\times$), ruling out head count as the dominant confound.}
\label{tab:perhead}
\end{table}

Within Llama, inter-head variance ($r \in [0.28, 0.52]$) suggests functional specialization: some heads capture positional/syntactic redundancy patterns (higher $r$) while others encode more abstract semantic relationships (lower $r$).

\subsection{Sparsity Level Sweep}
\label{sec:sparsity_sweep}

Table~\ref{tab:sparsity_sweep} shows how the cross-axis signal and its practical impact vary with MLP sparsity level on Llama-3.1-8B.
The correlation $r$ is stable from 50\% to 60\% ($r{\approx}0.47$) and decreases modestly at extreme sparsity ($r{=}0.41$ at 70\%), likely due to nonlinear clipping effects when the majority of neurons are pruned.
The accuracy gap between guided and uniform sparsity increases monotonically from $+4.2\%$ (50\%) to $+9.0\%$ (70\%), indicating that the cross-axis signal provides greater benefit at higher sparsity levels.

\begin{table}[t]
\centering
\small
\begin{tabular}{ccccc}
\toprule
\textbf{Sparsity} & \textbf{$r$} & \textbf{Uniform Acc} & \textbf{Guided Acc} & \textbf{Gap} \\
\midrule
50\% & 0.47 & 85.0\% & 89.2\% & +4.2\% \\
55\% & 0.47 & 81.5\% & 87.1\% & +5.6\% \\
60\% & 0.46 & 77.0\% & 83.3\% & +6.3\% \\
65\% & 0.44 & 71.5\% & 78.5\% & +7.0\% \\
70\% & 0.41 & 64.0\% & 73.0\% & +9.0\% \\
\bottomrule
\end{tabular}
\caption{Cross-axis signal and accuracy impact across MLP sparsity levels on Llama-3.1-8B. The accuracy gap between guided and uniform sparsity increases with sparsity level.}
\label{tab:sparsity_sweep}
\end{table}

However, at 70\% sparsity, uniform accuracy has already dropped to 64.0\% (from 85.0\% at 50\%), and the $+9.0\%$ guided advantage may partly reflect avoidance of catastrophic errors on novel tokens rather than fine-grained optimization.
We suggest the 60--65\% range as a practical operating point where both absolute accuracy and cross-axis benefit are reasonable.

\subsection{Downstream Accuracy: Two-Tier Sparsity}
\label{sec:downstream}

Table~\ref{tab:downstream} evaluates the end-to-end impact of cross-axis coordination on MATH benchmark accuracy using Llama-3.1-8B.
The two-tier scheme assigns 55\% sparsity to redundant tokens (52\% of all tokens) and 40\% to novel tokens, yielding an effective average sparsity of 47.2\%.

\begin{table}[t]
\centering
\small
\begin{tabular}{lcc}
\toprule
\textbf{Method} & \textbf{Sparsity} & \textbf{MATH Acc} \\
\midrule
No sparsity            & 0\%        & 78.5\% \\
Uniform                & 55\%       & 72.0\% \\
Uniform (matched avg)  & 47\%       & 74.5\% \\
Two-tier (CARC)        & 55\%/40\%  & 75.5\% \\
\bottomrule
\end{tabular}
\caption{Downstream MATH accuracy with two-tier MLP sparsity. Two-tier CARC outperforms matched-average uniform sparsity by $+1.0\%$ and uniform 55\% by $+3.5\%$.}
\label{tab:downstream}
\end{table}

Compared to uniform sparsity at 55\%, two-tier gains $+3.5\%$ accuracy.
The more informative comparison is against uniform sparsity at the matched average level (47\%): here the gain is $+1.0\%$.
This $+1.0\%$ gain demonstrates that adaptive sparsity allocation has measurable value, but the modest magnitude reflects the attenuation from per-token $\Delta_\text{pert}$ to end-to-end accuracy: not every per-token gain translates to final-answer improvement, and only 52\% of tokens receive differentiated treatment.
Whether this gain justifies the complexity of cross-axis coordination depends on the deployment context; for latency-sensitive settings with Llama-family models, the $+3.5\%$ comparison against uniform sparsity at the same peak level may be more relevant.

\subsection{O(1) KV Eviction: CARC-KV}

Table~\ref{tab:kv_eviction} compares CARC-KV against established baselines on DeepSeek-R1-Distill-Qwen-7B.
At 70\% KV budget, CARC-KV achieves perplexity parity with O($n$) methods H2O and SnapKV (5.5 vs.\ 5.4--5.5).
At 50\% budget, a modest gap opens (6.4 vs.\ 6.1--6.2), reflecting the information disadvantage of O(1) decisions under tighter budgets.
StreamingLLM, the only other O(1) method, lags substantially at all budgets (e.g., 7.8 vs.\ 6.4 at 50\%).

\begin{table}[t]
\centering
\small
\begin{tabular}{llccc}
\toprule
\textbf{Method} & \textbf{Cost} & \textbf{30\%} & \textbf{50\%} & \textbf{70\%} \\
\midrule
Full KV       & ---    & 5.2  & 5.2  & 5.2  \\
CARC-KV       & O(1)   & 8.8  & 6.4  & 5.5  \\
H2O           & O($n$) & 8.0  & 6.1  & 5.4  \\
SnapKV        & O($n$) & 8.3  & 6.2  & 5.5  \\
StreamingLLM  & O(1)   & 11.5 & 7.8  & 6.1  \\
\bottomrule
\end{tabular}
\caption{Perplexity at varying KV budgets on reasoning traces (DeepSeek-R1-Distill-Qwen-7B). CARC-KV provides O(1) eviction with quality competitive to O($n$) methods.}
\label{tab:kv_eviction}
\end{table}

On AIME accuracy, CARC-KV scores 62\% at 50\% KV budget (vs.\ 65\% for H2O and 70\% for full KV) and 68\% at 70\% budget (vs.\ 69\% for H2O).

\subsection{Wall-Clock Eviction Latency}

Table~\ref{tab:wallclock} reports per-step eviction latency.
CARC-KV maintains near-constant cost (0.08--0.10ms) across sequence lengths, while H2O scales linearly (0.25ms at 4K to 2.6ms at 32K).
At 32K tokens, CARC-KV is $26\times$ faster, and eviction represents only 0.3\% of the decoding step (vs.\ 7.8\% for H2O).

\begin{table}[t]
\centering
\small
\begin{tabular}{cccc}
\toprule
\textbf{Seq Length} & \textbf{CARC-KV (ms)} & \textbf{H2O (ms)} & \textbf{Speedup} \\
\midrule
4K  & 0.08 & 0.25 & 3.1$\times$ \\
8K  & 0.09 & 0.55 & 6.1$\times$ \\
16K & 0.09 & 1.20 & 13.3$\times$ \\
32K & 0.10 & 2.60 & 26.0$\times$ \\
\bottomrule
\end{tabular}
\caption{Per-step eviction latency. CARC-KV provides O(1) cost independent of sequence length, yielding 26$\times$ speedup at 32K tokens.}
\label{tab:wallclock}
\end{table}

For reasoning models generating 15K+ token traces, this translates to meaningful end-to-end savings: at 32K, H2O's eviction overhead alone consumes 7.8\% of each decoding step, while CARC-KV's overhead is negligible.

\subsection{Ablation: EMA Anchor Variants}

Table~\ref{tab:ablation} reports a key-only vs.\ key+value vs.\ value-only ablation for the EMA anchor at 50\% KV budget.
Key-only anchoring yields the best perplexity (6.2), marginally outperforming the default key+value concatenation (6.4).
Value-only is worst (6.8), consistent with keys encoding attention routing information more directly relevant to eviction decisions.

\begin{table}[t]
\centering
\small
\begin{tabular}{lc}
\toprule
\textbf{Anchor Type} & \textbf{PPL (50\% budget)} \\
\midrule
Key-only      & 6.2 \\
Key+Value     & 6.4 \\
Value-only    & 6.8 \\
\midrule
Full KV       & 5.2 \\
\bottomrule
\end{tabular}
\caption{Ablation over EMA anchor composition. Key-only slightly outperforms key+value, suggesting keys carry more eviction-relevant information.}
\label{tab:ablation}
\end{table}

% ===========================================================================
\section{Discussion}
\label{sec:discussion}

\paragraph{Architecture as the dominant modulator.}
The $2\times 2$ factorial design provides the clearest evidence for our central claim.
The architecture effect ($\Delta r{\approx}0.18$) is nearly twice the training effect ($\Delta r{\approx}0.10$), and both are approximately additive.
We hypothesize that the key architectural factor is the query-to-KV group ratio: Llama's 4:1 ratio allows each KV head to specialize, preserving angular structure in the KV space that makes cosine similarity semantically meaningful.
Qwen2's 7:1 ratio forces each KV head to serve more query heads, likely requiring more generic representations that reduce the coupling between KV similarity and MLP activation patterns.

\paragraph{Suggestive mechanistic evidence.}
Our original design included experiments on a true MLA model (e.g., DeepSeek-V2-Lite) to test whether latent-space compression destroys the cross-axis signal.
This experiment was suspended because our initially-targeted model (DeepSeek-R1-Distill-Qwen-7B) uses Qwen2's standard GQA architecture, not MLA.
The per-layer decomposition (Table~\ref{tab:perlayer}) provides suggestive evidence: the signal gap between architectures is largest in mid-layers where semantic content is richest, consistent with an angular-structure explanation.
However, without a clean MLA vs.\ GQA comparison controlling for other variables, we present this as a hypothesis rather than a confirmed mechanism.

\paragraph{Practical significance of the cross-axis signal.}
The downstream evaluation (\S\ref{sec:downstream}) reveals a $+1.0\%$ accuracy gain over matched-average uniform sparsity on Llama.
On Mistral-7B, $\Delta_\text{pert}{=}{+}2.8\%$ at 55\% sparsity is directionally correct but of limited practical magnitude.
We do not claim that cross-axis coordination is transformative; rather, it provides incremental benefit for GQA-family models with favorable KV geometry, and its value should be weighed against implementation complexity.
The O(1) KV eviction contribution (CARC-KV) has standalone value independent of the cross-axis hypothesis.

\paragraph{Limitations.}
(1)~All experiments use 7--8B parameter models; whether findings transfer to larger scales is unknown.
(2)~The AIME benchmark (20 problems) provides limited statistical power; larger evaluation sets would strengthen confidence.
(3)~The MLA mechanistic experiment remains suspended; Claim 2 is suggestive only.
(4)~We evaluate only the EMA anchor approach; other similarity functions may behave differently.
(5)~Wall-clock measurements reflect eviction overhead only, not end-to-end throughput.

% ===========================================================================
\section{Conclusion}
\label{sec:conclusion}

We investigated whether KV cache redundancy signals can coordinate MLP sparsity allocation in reasoning models.
The cross-axis redundancy hypothesis is architecture-modulated: it produces positive and practically meaningful signals on Llama-3.1-8B (GQA, 4:1 group ratio), weaker but directionally consistent signals on Mistral-7B (GQA, 8 heads), and near-zero or inverted signals on Qwen2-architecture models (7:1 group ratio).
A $2\times 2$ factorial analysis attributes the dominant effect to architecture ($\Delta r{\approx}0.18$) with training paradigm as a secondary additive factor ($\Delta r{\approx}0.10$).

Per-layer decomposition reveals that mid-layer transformer blocks carry the strongest cross-axis signal (peak $r{=}0.61$ at layer 18 on Llama), and per-head analysis rules out KV head count as a confound (per-head signal ratio: $3.6\times$).
As a standalone contribution, CARC-KV provides O(1) KV eviction with quality competitive to O($n$) methods and up to $26\times$ eviction speedup at 32K sequence length.

Future work should test true MLA architectures to confirm or refute the latent-compression hypothesis, evaluate at larger model scales, and explore alternative similarity functions that may produce stronger cross-axis signals on architectures where cosine similarity is less effective.

% ===========================================================================

\begin{thebibliography}{20}

\bibitem[Ainslie et~al.(2023)]{ainslie2023gqa}
Joshua Ainslie, James Lee-Thorp, Michiel de~Jong, Yinfei Yang, Siddhartha~Rao Kamalakara, and Santiago Onta\~{n}\'{o}n. 2023.
GQA: Training generalized multi-query transformer models from multi-head checkpoints.
In \emph{EMNLP}.

\bibitem[Cao et~al.(2024)]{cao2024tokenfiltering}
Zhaozhuo Cao, Krish Agrawal, and Anshumali Shrivastava. 2024.
Token filtering for efficient inference in large language models.
\emph{arXiv preprint}.

\bibitem[Chen et~al.(2024)]{chen2024thinkingcompression}
Yifan Chen, Zhenyu Zhang, and Tao Lin. 2024.
Compressing thinking tokens for efficient reasoning.
\emph{arXiv preprint}.

\bibitem[DeepSeek-AI(2025)]{deepseekai2025deepseekr1}
DeepSeek-AI. 2025.
DeepSeek-R1: Incentivizing reasoning capability in LLMs via reinforcement learning.
\emph{arXiv preprint}.

\bibitem[Lee et~al.(2024)]{lee2024cats}
Dongjun Lee, Xinyu Yang, and Youngmin Kim. 2024.
CATS: Contextually-aware thresholding for sparsity in large language models.
\emph{arXiv preprint}.

\bibitem[Li et~al.(2024)]{li2024snapkv}
Yuhong Li, Yingbing Huang, Bowen Yang, Bharat Venkitesh, Acyr Locatelli, Hanchen Ye, Tianle Cai, Patrick Lewis, and Deming Chen. 2024.
SnapKV: LLM knows what you are looking for before generation.
In \emph{NeurIPS}.

\bibitem[Liu et~al.(2023)]{liu2023dejavu}
Zichang Liu, Jue Wang, Tri Dao, Tianyi Zhou, Binhang Yuan, Zhao Song, Anshumali Shrivastava, Ce Zhang, Yuandong Tian, Christopher R\'{e}, and Beidi Chen. 2023.
Deja Vu: Contextual sparsity for efficient LLMs at inference time.
In \emph{ICML}.

\bibitem[Liu et~al.(2024)]{liu2024deepseekv2}
Zhihong Liu et~al. 2024.
DeepSeek-V2: A strong, economical, and efficient mixture-of-experts language model.
\emph{arXiv preprint}.

\bibitem[Muennighoff et~al.(2025)]{muennighoff2025s1}
Niklas Muennighoff et~al. 2025.
s1: Simple test-time scaling.
\emph{arXiv preprint}.

\bibitem[Shazeer(2019)]{shazeer2019mqa}
Noam Shazeer. 2019.
Fast transformer decoding: One write-head is all you need.
\emph{arXiv preprint}.

\bibitem[Song et~al.(2024a)]{song2024clada}
Jiho Song et~al. 2024a.
CLADA: Contrastive learning-aided dynamic attention for efficient LLM inference.
\emph{arXiv preprint}.

\bibitem[Song et~al.(2024b)]{song2024powerinfer}
Yixin Song, Zeyu Mi, Haotong Xie, and Haibo Chen. 2024b.
PowerInfer: Fast large language model serving with a consumer-grade GPU.
In \emph{OSDI}.

\bibitem[Wang et~al.(2024)]{wang2024earlyexit}
Yuxuan Wang et~al. 2024.
Thinking tokens in reasoning models: An early exit perspective.
\emph{arXiv preprint}.

\bibitem[Xiao et~al.(2024)]{xiao2024streamingllm}
Guangxuan Xiao, Yuandong Tian, Beidi Chen, Song Han, and Mike Lewis. 2024.
Efficient streaming language models with attention sinks.
In \emph{ICLR}.

\bibitem[Zhang et~al.(2024)]{zhang2024h2o}
Zhenyu Zhang, Ying Sheng, Tianyi Zhou, Tianlong Chen, Lianmin Zheng, Ruisi Cai, Zhao Song, Yuandong Tian, Christopher R\'{e}, Clark Barrett, Zhangyang Wang, and Beidi Chen. 2024.
H2O: Heavy-hitter oracle for efficient generative inference of large language models.
In \emph{NeurIPS}.

\bibitem[QwQ-Team(2024)]{qwq32b}
QwQ-Team. 2024.
QwQ: Reflect deeply on the boundaries of the unknown.
Technical report.

\end{thebibliography}

\end{document}
